\documentclass[11pt,a4paper]{article}

\usepackage[margin=2.5cm]{geometry}
\usepackage{hyperref}
\usepackage{amsmath}
\usepackage{booktabs}
\usepackage{enumitem}
\usepackage{titlesec}
\usepackage{xcolor}
\usepackage{parskip}
\usepackage{graphicx}
\usepackage{microtype}
\usepackage{multirow}

\hypersetup{
    colorlinks=true,
    linkcolor=blue!60!black,
    urlcolor=blue!60!black,
    citecolor=blue!60!black
}

\titleformat{\section}{\large\bfseries}{}{0em}{}[\titlerule]
\titleformat{\subsection}{\normalsize\bfseries}{}{0em}{}

\title{%
    \textbf{Generative Retrieval for Fine-Grained Chest X-Ray Case Retrieval}\\[0.4em]
    \large Research Plan and Progress Report\\[0.2em]
    \normalsize IRLab Internship --- University of Amsterdam
}
\author{Tea Cetojevic Tisaj \\ \small Supervised by Prof.\ Maarten de Rijke \& Yubao Tang}
\date{April - July 2026}

\begin{document}

\maketitle

\begin{abstract}
This report describes an internship project on \emph{generative multimodal
retrieval for fine-grained chest X-ray (CXR) case retrieval}.
Given a query in the form of a CXR image, a textual clinical description, or
both, the goal is to retrieve clinically relevant historical cases from a
study-level corpus.
The primary evaluation setting is \emph{non-overlapped case retrieval}: query
cases used for validation and testing are disjoint from the training corpus,
reflecting the realistic clinical scenario in which a new case is presented and
similar historical cases are sought.

The central research question is not simply whether the GENIUS generative
retrieval framework can be adapted to a new domain, but whether the paradigm
itself --- replacing nearest-neighbour search over continuous embeddings with
autoregressive decoding of discrete document identifiers --- can acquire
\emph{discriminative} semantic identifiers in a domain where images are globally
homogeneous and clinically meaningful variation is subtle.
We use BiomedCLIP as a domain-appropriate frozen encoder and retrain both the
residual quantizer and the T5 decoder on MIMIC-CXR.
The primary evaluation is in-domain (MIMIC-CXR test split); a secondary
evaluation assesses cross-dataset generalisation to IU~X-Ray.
If the standard pipeline achieves adequate discrimination, the result directly
validates generative retrieval for medical case search.
If not, the contribution shifts to targeted improvements to the quantization
objective, incorporating report-aware, disease-aware, or hard-negative-aware
supervision to encourage the codebook to separate clinically distinct cases.
This document describes the research framing, technical decisions, current
implementation status, and planned next steps.
\end{abstract}

\tableofcontents
\vspace{1em}

%% ---------------------------------------------------------------
\section{Motivation and Research Positioning}

Radiology image interpretation is inherently case-based: clinicians routinely
search prior studies with similar findings to support diagnostic reasoning,
triage, and treatment planning.
Existing retrieval systems overwhelmingly rely on \emph{embedding-based}
similarity search --- nearest-neighbour retrieval over dense continuous
embeddings --- returning cases ranked by cosine distance in some representation
space.
While effective, this approach requires maintaining and searching a vector index
that scales linearly with corpus size, and the retrieval objective (proximity
in embedding space) is not directly aligned with clinical relevance.

\emph{Generative retrieval} offers a structurally different paradigm.
Rather than representing candidates as points in a continuous vector space and
ranking by similarity, the system autoregressively \emph{generates} a discrete
semantic identifier for the most relevant candidate given a query.
This collapses retrieval to a sequence generation problem, offers near-constant
latency with respect to corpus size, and --- crucially --- embeds the mapping
from queries to cases directly into model parameters rather than into an
external index.

Chest X-rays present an especially demanding testbed for this paradigm.
Images are globally similar: all show lungs, ribs, and soft tissue in a largely
fixed anatomical geometry.
Yet clinically meaningful differences are subtle ---
a focal consolidation, an asymmetric pleural effusion, a widened mediastinum ---
and distinguishing these requires fine-grained discrimination rather than
coarse visual matching.
This creates a fundamental tension with the residual quantization step that
constructs discrete identifiers: codebook entries learned primarily from global
embedding statistics may fail to separate cases that a radiologist would
consider diagnostically distinct.

\textbf{Research positioning.}
This project does \emph{not} frame itself as a straightforward application of
GENIUS to a new domain.
Rather, the central question is whether the generative retrieval paradigm
itself --- specifically, the discrete identifier learning component --- can
achieve fine-grained discrimination in a clinically demanding,
low-visual-variance domain.
The answer shapes the primary contribution.
If the standard pipeline succeeds, this constitutes a positive result for
generative retrieval in medical case search.
If it fails --- identifiers collapse to coarse disease categories, codebook
utilisation degrades, or retrieval quality falls short of the dense baseline ---
the contribution lies in characterising the failure mode and proposing targeted
improvements to the quantization objective.

%% ---------------------------------------------------------------
\section{Background: GENIUS}

GENIUS~\cite{genius2025} is a generative retrieval framework for universal
multimodal search, achieving state-of-the-art performance among generative
retrieval methods on the M-BEIR benchmark (5.6M candidates across eight tasks).
The framework consists of three training stages:

\begin{enumerate}[itemsep=2pt]
    \item \textbf{Stage~0 --- Encoder.}
    CLIP ViT-L/14 is fine-tuned with contrastive learning on multimodal
    query-target pairs using natural-language instructions to align query
    intent with target semantics.

    \item \textbf{Stage~1 --- Residual Quantization (RQ).}
    Candidate embeddings are compressed into discrete ID sequences by a
    hierarchical residual quantizer.
    The first code captures modality (image, text, or image-text pair);
    subsequent codes encode semantics coarse-to-fine over a 4096-token
    codebook repeated across 8 levels.

    \item \textbf{Stage~2 --- Autoregressive Decoder.}
    A T5-small decoder is trained to generate the target candidate's ID
    sequence given a query embedding.
    At inference, trie-constrained beam search over the set of valid ID
    sequences prevents generation of identifiers that do not correspond to
    any corpus entry.
\end{enumerate}

GENIUS achieves nearly constant retrieval latency as corpus size grows,
unlike embedding-based methods whose search cost scales linearly with the
number of indexed candidates.

%% ---------------------------------------------------------------
\section{Problem Formulation}

\subsection{Project definition}

We define the project as \textbf{generative multimodal retrieval for
fine-grained chest X-ray case retrieval}.
Given a query in the form of a CXR image, a textual clinical description
(e.g.\ the impression or indication), or both, the system retrieves
\emph{clinically relevant historical cases} from a study-level corpus.
Each retrieved item is a study-level case that may contain one or more
radiograph images, free-text report sections, and structured pathology labels.

\textbf{Retrieval setting.}
The primary setting is \emph{non-overlapped case retrieval}: the cases used for
validation and testing do not appear in the training corpus.
This reflects the realistic clinical deployment scenario in which a new,
previously unseen case is presented to the system and similar historical cases
are retrieved to support diagnostic reasoning.
A \emph{fixed-corpus} setting, in which queries may self-match (i.e.\ a case
can retrieve itself), is a simpler configuration and may serve as an
exploratory comparison, but it is not the main evaluation target.

\subsection{Generative retrieval formulation}

Generative retrieval replaces nearest-neighbour search with autoregressive
decoding of a discrete \emph{semantic case identifier}.
Each study-level case $c$ is encoded into a fixed-length ID sequence
$T_c = (t^c_1, \ldots, t^c_M)$ via the following pipeline:

\[
\underbrace{(\text{image}_c,\; \text{report}_c)}_{\text{case}}
\;\xrightarrow{\text{medical encoder}}\;
\mathbf{e}_c \in \mathbb{R}^d
\;\xrightarrow{\text{quantizer}}\;
T_c \in \mathcal{V}^M
\]

where $\mathcal{V}$ is the codebook vocabulary (4096 tokens per level) and
$M = 8$ semantic levels are used, following the GENIUS configuration.

Given a query $q$ (image, text, or both), the retriever generates the
identifier of the most relevant case autoregressively:
\[
t^c_k = \arg\max_{t \in \mathcal{V}} \left[
    \log p\!\left(t \;\middle|\; \mathbf{q},\, t^c_{<k};\, \theta\right)
\right],
\]
where $\mathbf{q}$ is the query embedding produced by the same medical encoder.
Generation is constrained by a trie over valid identifier sequences, ensuring
only IDs that correspond to corpus entries can be produced.

A key property of this formulation is that the semantic ID is directly the case
label: since each identifier corresponds to exactly one study, retrieval
evaluation can be performed directly on the generated ID without a separate
re-ranking or index lookup step.

\subsection{Candidate representation}

Each candidate is defined at the \textbf{study level}:
\begin{itemize}[itemsep=2pt]
    \item \textbf{Image}: frontal-view radiograph (PA preferred; AP as fallback).
    Studies with no frontal view are excluded.
    \item \textbf{Text}: concatenation of the \textit{Findings} and
    \textit{Impression} sections of the associated radiology report,
    following standard practice in the CXR literature.
    \item \textbf{Candidate modality}: \texttt{image,text} (multimodal).
    \item \textbf{Query modality}: \texttt{image}, \texttt{text}, or
    \texttt{image,text}, depending on the query type.
    The initial evaluation uses image-only queries, reflecting the realistic
    clinical scenario in which a clinician searches by radiograph alone.
\end{itemize}

Multimodal candidates (\texttt{image,text}) are used rather than text-only or
image-only: the combined representation provides richer quantizer supervision
and preserves all information available at indexing time.
The M-BEIR task type for image query against multimodal candidates is
\texttt{image$\to$image,text} (task~9), which was absent from the original
\texttt{MBEIR\_TASK} mapping in \texttt{utils.py} and has been added.
Note that this is distinct from task~7 (\texttt{image,text$\to$image}), which
requires a multimodal query.

\subsection{Research questions}

\textbf{RQ1 (primary):}
Can the generative retrieval paradigm, adapted to the medical domain via a
frozen biomedical vision-language encoder, learn sufficiently discriminative
semantic case identifiers for fine-grained CXR case retrieval under the
non-overlapped setting?

\textbf{RQ2:}
How does generative retrieval compare to an embedding-based dense retrieval
baseline using the same BiomedCLIP encoder?
This comparison isolates the contribution of the generative framework from
that of the encoder.

\textbf{RQ3:}
What are the failure modes of residual quantization under visual homogeneity?
Does codebook collapse occur, and at what level of the quantization hierarchy
does discriminability degrade?

\textbf{RQ4:}
If the standard pipeline is insufficient, which training signal ---
report-aware, disease-aware, or hard-negative-aware supervision --- most
effectively improves identifier discriminability?

\textbf{RQ5:}
Does the trained system generalise across datasets
(MIMIC-CXR $\to$ IU~X-Ray), or does it overfit to MIMIC-CXR-specific
acquisition and reporting conventions?

%% ---------------------------------------------------------------
\section{Proposed Approach}

\subsection{Encoder: BiomedCLIP (frozen)}

We replace CLIP ViT-L/14 with BiomedCLIP~\cite{biomedclip2024}
(ViT-B/16, 512-dimensional output), pretrained on 15~million biomedical
image-text pairs from PubMed Central.
BiomedCLIP has been shown to outperform general-purpose CLIP on a range of
medical imaging benchmarks and provides a domain-appropriate representation
space for chest radiographs.

The encoder is used \emph{frozen}: its weights are not updated during
Stages~1 or~2.
Stage~0 fine-tuning of BiomedCLIP is deferred as a potential extension;
fine-tuning under a contrastive objective would require careful design of
hard negatives for CXR (see Section~\ref{sec:open}) and is not necessary to
evaluate the core retrieval framework.

The change from 768-dimensional CLIP embeddings to 512-dimensional BiomedCLIP
embeddings is absorbed directly: the quantizer input dimension and the T5
decoder projection are set to 512 in the configuration, with no intermediate
adapter.
Both Stage~1 and Stage~2 are retrained from scratch using BiomedCLIP features.

\subsection{Dataset: MIMIC-CXR}

MIMIC-CXR-JPG~\cite{mimicjpg2019} contains approximately 370,000 radiograph
images from around 227,000 studies, provided with structured
train/validation/test splits that are stratified at the patient level to
prevent data leakage.

\begin{itemize}[itemsep=2pt]
    \item \textbf{Candidate pool}: all train-split studies with a frontal view
    and a non-empty report ($\approx$180--200K studies).
    \item \textbf{Stage~1 training}: BiomedCLIP embeddings of all train-split
    candidates.
    \item \textbf{Stage~2 training}: image queries paired with their
    corresponding study IDs as targets, using the train split.
    \item \textbf{Evaluation}: MIMIC-CXR test split (closed-corpus,
    in-distribution).
    \item \textbf{Cross-dataset evaluation}: IU~X-Ray ($\approx$3,955 studies),
    zero-shot or with fine-tuning on the IU~X-Ray training split.
\end{itemize}

\subsection{Training stages (adapted)}
\label{sec:stages}

\textbf{Stage~0} is \emph{skipped} at this stage.
BiomedCLIP is used as a pretrained frozen encoder.
Fine-tuning is listed as a future extension.

\textbf{Stage~1} (residual quantizer) is \emph{retrained from scratch} on
MIMIC-CXR train embeddings.
The codebook size (4096 tokens per level) and number of levels (8 semantic
levels plus 1 modality indicator) follow the default GENIUS configuration,
which serves as the baseline.
Codebook utilisation --- the fraction of active codes and entropy of the
assignment distribution --- will be monitored throughout training to detect
collapse.

\textbf{Stage~2} (T5 decoder) is \emph{retrained} on MIMIC-CXR train queries.
T5-small is initialised from pretrained weights (hidden dimension 512 matches
our embedding dimension); the output projection over the code-token vocabulary
is reinitialised randomly.

\subsection{Evaluation protocol}
\label{sec:eval}

\textbf{Retrieval setting.}
All evaluations use the non-overlapped setting: queries are drawn from the
MIMIC-CXR test split, which is patient-disjoint from the training corpus.
A query case cannot appear in the candidate pool, so self-match is not possible.
This mirrors the intended clinical deployment scenario and is the harder,
more realistic evaluation condition.

\textbf{Ground truth definition.}
Retrieval positives are defined by shared CheXpert pathology labels: a
candidate is considered relevant to a given query if the two studies share at
least one active (positive) pathology label as assigned by the CheXpert
labeller.
This definition is clinically interpretable, stable across dataset splits, and
requires no additional annotation beyond the labels distributed with
MIMIC-CXR-JPG.
Its main limitation --- insensitivity to severity and laterality --- is
acknowledged and treated as a direction for future refinement.

\textbf{Metrics.}
Following the M-BEIR convention: Recall@5 and Recall@10.
For cross-dataset generalisation (IU~X-Ray), the same metrics apply.

\textbf{Baseline.}
The primary comparison is an embedding-based dense retrieval system using the
same frozen BiomedCLIP encoder: query and candidate embeddings are compared by
cosine similarity with no learned quantization or autoregressive decoder.
This comparison isolates the contribution of the generative retrieval framework
from the representation quality of the encoder itself.

%% ---------------------------------------------------------------
\section{Technical Implementation}

\subsection{Codebase modifications}

The GENIUS codebase required targeted modifications to support a different
encoder and embedding dimension.
Table~\ref{tab:impl} summarises the components implemented to date.

\begin{table}[h]
\centering
\small
\caption{Implementation status of pipeline components.}
\label{tab:impl}
\begin{tabular}{@{}lll@{}}
\toprule
\textbf{Component} & \textbf{File} & \textbf{Status} \\
\midrule
BiomedCLIP encoder wrapper & \texttt{src/models/biomedclip/biomedclip\_nofusion.py} & Done \\
Feature extraction script & \texttt{src/feature\_extraction/biomedclip\_feature\_extraction.py} & Done \\
Feature extraction config & \texttt{src/feature\_extraction/config\_biomedclip\_cand.yaml} & Done \\
Stage~1 RQ config (\texttt{feature\_dim: 512}) & \texttt{configs\_scripts/.../inbatch\_biomedclip.yaml} & Done \\
Stage~2 generator config & \texttt{configs\_scripts/.../inbatch\_biomedclip.yaml} & Done \\
Configurable \texttt{feature\_dim} (768 $\to$ 512) & \texttt{residual\_quantization.py}, \texttt{retriever.py} & Done \\
MIMIC-CXR preprocessing script & \texttt{src/data/preprocessing/mimic\_cxr.py} & Done \\
Dataset registration in \texttt{utils.py} & \texttt{DATASET\_IDS}, \texttt{MBEIR\_DATASET\_TO\_DOMAIN} & Done \\
Eval config for MIMIC-CXR & \texttt{config\_eval\_mimic\_cxr.yaml} & Done \\
Slurm job scripts & \texttt{scripts/run\_biomedclip\_*.sh} & Done \\
IU~X-Ray preprocessing (smoke test) & \texttt{src/data/preprocessing/iu\_xray.py} & Done \\
\texttt{open\_clip\_torch} dependency & cluster \texttt{genius2} environment & Installed \\
\bottomrule
\end{tabular}
\end{table}

\subsection{BiomedCLIP encoder wrapper}

The wrapper (\texttt{biomedclip\_nofusion.py}) provides the same interface as
the existing \texttt{CLIPNoFusion} class used in GENIUS, so the rest of the
pipeline requires no changes to its encoder-facing API.
The key method is \texttt{encode\_multimodal\_input(img, txt)}, which returns a
pair of \texttt{[B, 512]} tensors for image and text embeddings respectively.
The wrapper handles the HuggingFace tokenizer output format required by
\texttt{open\_clip}, which differs from the standard CLIP tokenizer interface.
Encoder weights are frozen throughout; no gradients flow into BiomedCLIP.

\subsection{MIMIC-CXR preprocessing}

The preprocessing script (\texttt{mimic\_cxr.py}) operates on the MIMIC-CXR-JPG
directory structure and produces M-BEIR-compatible JSONL files for both the
candidate pool and query sets.
For each study, it selects the frontal-view image (PA preferred over AP),
extracts the Findings and Impression sections via case-insensitive regex (to
handle header variants), and constructs a study-level candidate entry.
The CheXpert pathology labels are stored in the \texttt{src\_content} field of
each candidate entry --- not used during training but available for
disease-aware supervision in future experiments.
Studies with no frontal view or with report text shorter than ten characters
are excluded with a logged warning.
Train/validation/test assignment follows the official MIMIC-CXR split CSV,
which is patient-stratified.

\subsection{IU X-Ray as end-to-end smoke test}

Because MIMIC-CXR access requires PhysioNet credentialing (currently in
progress), IU~X-Ray~\cite{iuxray2015} ($\approx$7,500 studies, publicly available)
is used as a smoke-test corpus.
The goal is to verify that preprocessing, feature extraction, Stage~1, Stage~2,
and evaluation all run to completion without errors \emph{before} MIMIC-CXR
data arrives.
No meaningful retrieval results are expected at this scale, and codebook
collapse is anticipated; the exercise is purely diagnostic.

%% ---------------------------------------------------------------
\section{Technical Challenges}

\begin{enumerate}[itemsep=4pt]

    \item \textbf{Quantizer discriminability under visual homogeneity.}
    The residual quantizer is trained to minimise reconstruction error in
    embedding space.
    If BiomedCLIP embeddings for clinically distinct CXR studies are close in
    Euclidean distance --- plausible given global image similarity ---
    the quantizer may assign identical ID sequences to cases a radiologist
    would distinguish.
    Codebook utilisation and per-level assignment entropy will be monitored
    during Stage~1 training.

    \item \textbf{Corpus scale.}
    GENIUS was trained on 1.5M+ queries and 5.6M candidates; MIMIC-CXR provides
    approximately 200K training studies.
    At this scale, the codebook may be over-parameterised relative to the
    training data, risking dead codes (unused codebook entries) or uneven
    utilisation that degrades retrieval diversity.

    \item \textbf{In-batch negative quality.}
    GENIUS Stage~2 training uses in-batch negatives.
    In a general-domain corpus, randomly sampled negatives are typically
    easy to distinguish from the true positive.
    For chest X-rays, in-batch negatives may be visually and semantically similar
    to the query (e.g.\ two studies both showing ``No finding''), providing
    a weak training signal.
    CheXpert labels stored in \texttt{src\_content} are retained to enable
    disease-aware hard negative mining without re-running preprocessing.

    \item \textbf{Multi-view fusion.}
    MIMIC-CXR studies typically include a frontal and a lateral view.
    The current approach uses only the frontal view; whether multi-view
    fusion (e.g.\ mean-pooled embedding over available views) improves
    representation quality for the quantizer is left as a future experiment.

    \item \textbf{Ground truth granularity.}
    CheXpert-label-overlap as the ground truth definition is coarse:
    two studies sharing the label ``Consolidation'' may represent very
    different clinical presentations.
    This is a known limitation that constrains the interpretability of
    Recall@K metrics but does not invalidate the relative comparison
    between the generative and embedding-based baselines.

\end{enumerate}

%% ---------------------------------------------------------------
\section{Progress Summary}

\begin{itemize}[itemsep=4pt]
    \item GENIUS codebase understood at the level of all three training stages
    and the inference pipeline.
    \item Conda environment \texttt{genius2} configured on the ILPS cluster
    (Python~3.10, PyTorch, \texttt{open\_clip\_torch} installed).
    \item All original GENIUS checkpoints downloaded:
    \texttt{clip\_sf\_large.pth}, \texttt{rq\_clip\_large.pth},
    \texttt{GENIUS\_t5small.pth}.
    \item COCO M-BEIR preprocessing complete and verified;
    JSONL candidate pool and query files generated.
    \item BiomedCLIP encoder wrapper implemented with GENIUS-compatible
    interface; feature extraction scripts and configs ready.
    \item \texttt{feature\_dim} made configurable throughout
    \texttt{residual\_quantization.py} and \texttt{retriever.py};
    all 768-hardcoded occurrences replaced.
    \item MIMIC-CXR preprocessing script written (ready to run on data arrival);
    dataset registered in \texttt{utils.py}.
    \item Stage~1 and Stage~2 configs for BiomedCLIP (\texttt{feature\_dim: 512})
    prepared.
    \item Eval config for MIMIC-CXR prepared (\texttt{config\_eval\_mimic\_cxr.yaml}).
    \item Slurm job scripts written for feature extraction, Stage~1, and Stage~2.
    \item IU~X-Ray preprocessing script written for smoke-test pipeline.
    \item PhysioNet credentialing application for MIMIC-CXR in progress.
    \item Literature review ongoing: GENIUS, GIT-CXR, FAST-MRG, IRGen, UniIR,
    BiomedCLIP.
\end{itemize}

%% ---------------------------------------------------------------
\section{Planned Next Steps}

\begin{enumerate}[itemsep=4pt]

    \item \textbf{IU~X-Ray smoke test.}
    Run the full pipeline (preprocessing $\to$ BiomedCLIP feature extraction
    $\to$ Stage~1 $\to$ Stage~2 $\to$ eval) on IU~X-Ray to verify end-to-end
    correctness before MIMIC-CXR data arrives.

    \item \textbf{Obtain and stage MIMIC-CXR data.}
    Complete PhysioNet credentialing; download MIMIC-CXR-JPG ($\approx$100--150~GB)
    to the cluster shared filesystem.

    \item \textbf{BiomedCLIP feature extraction on MIMIC-CXR.}
    Run \texttt{run\_biomedclip\_feature\_extraction.sh} on the MIMIC-CXR train
    and candidate splits (GPU Slurm job, 4$\times$ A6000).

    \item \textbf{Stage~1: quantizer training.}
    Train the residual quantizer on MIMIC-CXR train embeddings; log codebook
    utilisation, dead-code fraction, and per-level entropy.
    Tune codebook size and number of levels if collapse is observed.

    \item \textbf{Stage~2: decoder training.}
    Train the T5-small decoder on MIMIC-CXR train queries.

    \item \textbf{Evaluation (in-domain).}
    Evaluate on MIMIC-CXR test split (Recall@5, Recall@10).
    Compare against the BiomedCLIP embedding baseline.

    \item \textbf{Diagnosis of failure modes.}
    If performance is below the embedding baseline, analyse codebook assignment
    patterns: do semantically similar studies share identical prefixes?
    At which level does discrimination break down?

    \item \textbf{Quantizer improvements (conditional on diagnosis).}
    Depending on failure mode analysis, implement one or more of:
    disease-aware contrastive loss using CheXpert labels;
    report-similarity-weighted hard negatives (BioSentVec);
    or modified RQ training with pathology-conditioned codebook initialisation.

    \item \textbf{Cross-dataset evaluation.}
    Zero-shot evaluation on IU~X-Ray, or fine-tuning on the IU~X-Ray train
    split if zero-shot performance is insufficient to draw conclusions.

\end{enumerate}

%% ---------------------------------------------------------------
\section{Open Questions}
\label{sec:open}

\begin{enumerate}[itemsep=6pt]

    \item \textbf{Hard negative design.}
    In-batch negatives may be insufficient for fine-grained CXR discrimination.
    If so, which supervision signal is most effective: pathology-label-based
    hard negatives (same modality, different pathology), report-similarity-based
    negatives (high BERTScore, different patient), or distance-based negatives
    in BiomedCLIP embedding space?
    The choice affects both Stage~1 and Stage~2 training.

    \item \textbf{Stage~0 fine-tuning of BiomedCLIP.}
    The current plan uses BiomedCLIP frozen.
    If Stage~1 codebook quality is poor --- indicated by high dead-code fraction
    or low utilisation entropy --- fine-tuning the encoder under a
    contrastive objective on MIMIC-CXR report-image pairs may improve the
    embedding space prior to quantization.
    This is a resource-intensive extension and is deferred pending Stage~1
    diagnostic results.

    \item \textbf{Multi-view fusion.}
    MIMIC-CXR studies typically include frontal and lateral projections.
    The current approach selects only the frontal view.
    Whether mean-pooled multi-view embeddings improve codebook quality or
    retrieval performance is empirically unexplored in this setting.

\end{enumerate}

%% ---------------------------------------------------------------
\section{Target Venues}

\begin{itemize}[itemsep=4pt]
    \item \textbf{ECIR 2027} (primary) --- Information retrieval community;
    deadline approximately October--November 2026.
    IRLab UvA has a strong presence at ECIR and the venue is directly
    appropriate for retrieval methodology work.
    \item \textbf{SIGIR 2027} (backup) --- Deadline approximately
    January 2027; higher competitive bar; appropriate if results are strong.
    \item \textbf{ImageCLEF 2026} (if results arrive early) --- Medical image
    retrieval track; timeline to be confirmed with supervisor.
\end{itemize}

%% ---------------------------------------------------------------
\section*{References}
\begin{thebibliography}{9}

\bibitem{genius2025}
S.\ Kim, X.\ Zhu, X.\ Lin, M.\ Bastan, D.\ Gray, and S.\ Kwak.
\textit{GENIUS: A Generative Framework for Universal Multimodal Search.}
arXiv:2503.19868, 2025.

\bibitem{biomedclip2024}
S.\ Zhang et al.
\textit{BiomedCLIP: a multimodal biomedical foundation model pretrained from
fifteen million scientific image-text pairs.}
arXiv:2303.00915, 2023.

\bibitem{mimicjpg2019}
A.\ E.\ Johnson et al.
\textit{MIMIC-CXR-JPG, a large publicly available database of labeled
chest radiographs.}
arXiv:1901.07042, 2019.

\bibitem{iuxray2015}
D.\ Demner-Fushman et al.
\textit{Preparing a collection of radiology examinations for distribution
and retrieval.}
\textit{Journal of the American Medical Informatics Association},
23(2):304--310, 2016.

\end{thebibliography}

\end{document}
